\documentclass[11pt,reqno]{amsart}

\pdfinfoomitdate=1
\pdftrailerid{}
\pdfsuppressptexinfo=-1

\usepackage{microtype}
\microtypesetup{expansion=false}
\usepackage{mathtools}
\usepackage{amssymb}
\usepackage{url}
\usepackage[hidelinks]{hyperref}
\hypersetup{
  pdfauthor={},
  pdftitle={},
  pdfsubject={},
  pdfkeywords={},
  pdfcreator={},
  pdfproducer={}
}

\numberwithin{equation}{section}
\newtheorem{theorem}{Theorem}[section]
\newtheorem{lemma}[theorem]{Lemma}

\setlength{\emergencystretch}{1em}

\makeatletter
\newcommand{\doi}[1]{%
  \begingroup
  \def\UrlBreaks{\do\/\do\.\do\-\do\_\do\(\do\)}%
  \def\UrlBigBreaks{\do\/}%
  \href{https://doi.org/#1}{\nolinkurl{#1}}%
\endgroup}
\makeatother

\date{August 9, 2026}

\begin{document}

\section{Proof}

Let $\mathcal D$ be the top-aligned $5$-by-$5$ Ferrers diagram with column
heights $(c_1,\ldots,c_5)=(3,3,5,5,5)$.  Write
$\kappa_F(\mathcal D,5)$ for the largest dimension of an $F$-linear space
of matrices supported on $\mathcal D$ in which every nonzero matrix has
rank $5$.  The Etzion--Silberstein bound is
\begin{equation}
 \nu_{\min}(\mathcal D,5)
 =\min_{0\le j<5}\sum_{i=1}^{5-j}
   \max\{0,c_i-4+j\}.
 \label{eq:bound}
\end{equation}

\begin{theorem}[Etzion--Silberstein counterexample]
For $F=\mathbb F_{169}$,
\[
 \kappa_F(\mathcal D,5)=2<3=\nu_{\min}(\mathcal D,5).
\]
Consequently, the Etzion--Silberstein bound is not universally attained.
\end{theorem}

\begin{proof}
The five quantities minimized in \eqref{eq:bound} are
\[
 (3,4,5,4,3),
\]
so the bound is $3$.

We first prove that the optimum is at most $2$.  Put $q=169$ and suppose
that a three-dimensional code exists.  Choose a basis $M_0,M_1,M_2$, put
$M(x)=\sum_{i=0}^2x_iM_i$, and identify $x$ with a column vector.  The
Ferrers support gives
\[
 M(x)=\begin{pmatrix}P(x)&B(x)\\0&D(x)\end{pmatrix},
 \qquad P(x)=[A_1x,A_2x],
\]
where $A_1,A_2\in M_3(F)$.  Each $A_i$ is invertible, since otherwise
some nonzero $M(x)$ has a zero column.  Moreover,
$T=A_1^{-1}A_2$ has no eigenvalue in $F$, since an eigenvector would make
the first two columns of $M(x)$ dependent.  Thus the cubic characteristic
polynomial of $T$ is irreducible.  Over $E=\mathbb F_{q^3}$ it has an
eigenvector $v\ne0$, and hence
\begin{equation}
 \det M(v)=0.
 \label{eq:extension-zero}
\end{equation}

Now set $f(X)=\det(\sum_{i=0}^2X_iM_i)$.  This is a homogeneous quintic
with no nonzero zero over $F$.  It is irreducible over $F$: otherwise one
of its homogeneous factors has degree at most $2$, and such a form in
three variables has a nonzero zero over a finite field, by elementary
linear algebra in degree $1$ and Chevalley--Warning in degree $2$.

If $f$ were absolutely irreducible, the normalization of the plane curve
$f=0$ would have genus at most $6$.  Hasse--Weil would give at least
$169+1-12\cdot13=14$ rational points; their images
would be projective $F$-zeros of $f$, a contradiction.  Hence $f$ is not
absolutely irreducible.

Since $\operatorname{char}F=13>5$, the quintic is geometrically reduced.
Frobenius acts transitively on its absolutely irreducible factors, whose
number times their common degree is $5$.  Irreducibility over $F$ and
failure of absolute irreducibility therefore force five conjugate lines.
Thus, for $L=\mathbb F_{q^5}$ and a nonzero $c\in F$,
\begin{equation}
 f=c\prod_{i=0}^4\ell^{(q^i)},
 \qquad
 \ell=a_0X_0+a_1X_1+a_2X_2,\quad a_i\in L.
 \label{eq:norm}
\end{equation}
The coefficients $a_0,a_1,a_2$ are $F$-linearly independent.  Otherwise
$\ell$, and therefore $f$, would have a nonzero $F$-zero.  Because
$E\cap L=F$ (the extension degrees $3$ and $5$ are coprime), they remain
independent over $E$.  The same holds for their Frobenius conjugates, so
\eqref{eq:norm} has no nonzero zero in $E^3$.  This contradicts
\eqref{eq:extension-zero}.  No three-dimensional code exists, and a code
of larger dimension would contain one; hence
$\kappa_F(\mathcal D,5)\le2$.

For the reverse inequality, let $K=\mathbb F_{q^5}=F(\alpha)$ and use the
power basis $1,\alpha,\ldots,\alpha^4$.  For $a,b\in F$, let
$M_{a+b\alpha}$ be the matrix of multiplication by $a+b\alpha$ on $K$.
The space
\[
 C_0=\{M_{a+b\alpha}:a,b\in F\}
\]
has dimension $2$.  Its first two columns represent
$a+b\alpha$ and $a\alpha+b\alpha^2$, so their last two entries vanish;
there is no restriction on the remaining columns.  Thus $C_0$ is
supported on $\mathcal D$.  Every nonzero member is invertible because it
is multiplication by a nonzero field element.  Therefore
$\kappa_F(\mathcal D,5)\ge2$, completing the proof.
\end{proof}

\end{document}
